\documentclass{article}
\usepackage{graphicx} % Required for inserting images

\title{Facial Biometric Verification System - IRS}
\author{Banele Mdluli}
\date{February 2026}

\begin{document}

\maketitle
\newpage
\section{Summary}
Facial biometric verification is a technology that uses facial recognition algorithms to verify the identity of an individual based on their facial features. This technology has gained popularity in recent years due to its convenience and security benefits. The facial biometric verification system can be used in various applications, including access control, identity verification, and fraud prevention.

\section{Introduction}
Authentications and verifications are essential in various industries, including banking, healthcare, and government. Traditional methods of authentication, such as passwords and PINs, are vulnerable to hacking and identity theft. 
Facial biometric verification provides a more secure and convenient alternative to traditional methods. The system uses advanced algorithms to analyze facial features and compare them to a database of known faces to verify the identity of an individual.


\section{Problem statement}
The current methods are vulnerable to hackers, to improve the authentication checks, the project to aims to develop facial biometric system that can be integrated with other applications.


\section{Objectives}
The project objectives are targets that must be reached for the project to be considered a success. The objectives are:
\begin{itemize}
    \item Develop a facial biometric system to authenticate users based on their ID number and photo.
    \item Develop a system that integrates with other applications for general use and enhanced validation process.
\end{itemize}

\section{Requirements}
\subsection{Functional requirements}
\begin{itemize}
 \item Facial biometric: 
    \begin{itemize}
        \item Facial biometric initiation
        \begin{itemize}
            \item A user must select or press a buttton to start the facial biometric process.
            \item An information page should be displayed to a user briefly explaining the process and the information required.
            \item Information that is required should be an ID number, a device with a camera and environment with lighting.
        \end{itemize}
        \item Facial biometric link 
        \begin{itemize}
            \item Once a user proceeds from the initiation phase, a link should be generated that is sent to the user. The link should be sent via email or SMS.
            \item The link that is generated should expire with 3 minutes. 
            \item When the link is clicked, it should direct the user to the facial biometric site on the web.
        \end{itemize}
        \item Facial image capturing process 
        \begin{itemize}
            \item The system should inform the user of any user hardware or environment negatively affecting the image capturing.
            \item The system should capture multiple images that can be used for validation.
            \item The system should request a re-try when the captured images do not meet the image standards
        \end{itemize}
        \item Facial image processing
        \begin{itemize}
            \item The system should inform the customer that the images have been captured and validation is underway.
            \item The system should re-direct the user to the home-page and an indicator show that an validation is ongoing.
            \item The system should inform the customer when the validation is complete via Email and SMS. 
            Depending on the outcome, the system should send a successful or failed authentication to internal systems. 
            \item The system should not allow a user reperform facial biometric with an hour from the last attempt. 
        \end{itemize}
    \end{itemize}
    \item Facial biometric audit and tracing 
    \begin{itemize}
        \item The system should record transaction event using date and time.
        \item The system should generated a report for a single transaction and multiple transactions within a specific period. 
    \end{itemize}
\end{itemize}
\subsection{Non functional requirements}
\begin{itemize}
    \item Performance-based requirements
    \begin{itemize}
        \item The API should provide feedback within 5-7 secs. Or provide reasons for the delay.
    \end{itemize}
    \item Scalability requirements 
    \begin{itemize}
        \item The application should be packaged using containers and Kubernetes.
        \item During high peak traffic, the application should increase the number of running instances (VM).
        \item During low trough traffic, the application should decrease the number of running instances (VM).
    \end{itemize}
    \item Reliability requirements 
    \begin{itemize}
        \item The application should automatically restart itself when it encounters an issue that leads to downtime.
        \item The application should meet an 80.00\% uptime (calculated daily).
    \end{itemize}
    \item Maintainability requirements
    \begin{itemize}
        \item The application must have a production, development and testing(pre-production) environment.
        \item The application should use version control and a semantic versioning process.
        \item the application should have automated tests using continuous deployment pipeline.
    \end{itemize}
    \end{itemize}
\subsection{Technical requirement}
    \begin{itemize}
        \item A.I agent framework (\textbf{To be confirmed})
        \item API dev framework
        \begin{itemize}
            \item Fast API and documentation OpenAPI (Swagger)
        \end{itemize}
        \item Frontend 
        \begin{itemize}
            \item React
            \item Streamlit
        \end{itemize}
        \item Environment packaging 
        \begin{itemize}
            \item docker
            \item Kubernetes
        \end{itemize}
        \item Deployment platform, code repo and CI/CD
        \begin{itemize}
            \item Azure
            \item Github
            \item Github actions  (switch to Azure DevOps later)
    \end{itemize}
 \end{itemize}

 \section{User acceptance criteria}
 \begin{table}[h]
    \caption{User story acceptance criteria}
    \centering
    \begin{tabular}{|p{3cm}|p{5cm}|p{5cm}|}
        \hline
        \textbf{Acceptance Criterion} & \textbf{Given} & \textbf{When} \\
        \hline
        The system should provide a link to the user to perform facial biometric & 
        A user has selected to perform facial biometric &
         The system should generate a link and send it to the user via email or SMS.\\
        \hline
        The system should capture multiple images of the user for validation &
        A user has clicked the link and is on the facial biometric site &
        The system should capture multiple images of the user and inform the user of any hardware or environment negatively affecting the image capturing.\\
        \hline
        The system should inform the user of the validation outcome &
        The system has completed the validation process &
        The system should inform the user of the validation outcome via email or SMS. Depending on the outcome, the system should send a successful or failed authentication to internal systems.\\
        \hline
        The system should not allow a user to reperform facial biometric within an hour from the last attempt &
        A user has attempted to perform facial biometric within an hour from the last attempt &
        The system should inform the user that they cannot reperform facial biometric within an hour from the last attempt.\\
        \hline
        The system should record transaction events using date and time &
        A user has performed facial biometric &
        The system should record transaction events using date and time and generate a report for a single transaction and multiple transactions within a specific period.\\
        \hline 
        The system should provide feedback within 5-7 seconds or provide reasons for the delay &
        A user has performed facial biometric &
        The system should provide feedback within 5-7 seconds or provide reasons for the delay.\\
        \hline
        The application should automatically restart itself when it encounters an issue that leads to downtime &
        The application has encountered an issue that leads to downtime &
        The application should automatically restart itself when it encounters an issue that leads to downtime.\\
        \hline
        The application should meet an 80.00\% uptime (calculated daily) &
        The application is running &
        The application should meet an 80.00\% uptime (calculated daily).\\
        \hline
        The application should have a production, development and testing(pre-production) environment &
        The application is being developed &
        The application should have a production, development and testing(pre-production) environment.\\
        \hline
    \end{tabular}
 \end{table}

 \section{Design}
The low-level design shows the structure and make-up of individual units. The LLD is a platform divided; units generated or developed on the same platform are part of the same LLD design.\\

\section{High technical requirements}
\subsection{Fast API integrations}
\subsubsection{Postgres, Agent and Fast API}
 A PostgreSQL database should be used to store information analysed in Databricks. The information stored in the database will vary from data accuracy to validity checks.
 \begin{itemize}
     \item  Create an Azure Postgres table to store data quality outcomes. Depending on the outcome size, one table can be used, or multiple tables (one check per table: if required).
     \item  Create a Fast API parent endpoint that will interact with an agent or agents that will create a data quality check PDF report.
     \item Create a Decision matrix that should assist the agent in compiling the report to give a certain outcome.
 \end{itemize}
 \subsubsection{Gradio, React and Fast API}
    \begin{itemize}
        \item Create an agent that should interact with an external user and provide the desired outcome based on the information available on the PostgreSQL database. The agent should
        interact with the database using Fast API.
        \item Create a Gradio chat interface that is connected to the AI agent. Link the Gradio chat app to a React Interface using Fast API.
        \item Create a React platform that has the following:
            \begin{itemize}
                \item Log in page with an option to sign up.
                \item A page that shows operation indicators.
                \item A page that shows KPI based performance indicators.
                \item A page that allows users to interact with agents.
            \end{itemize}
        \item Report data should be downloadable in various formats (csv, tsv, etc,...).
    \end{itemize}
\subsection{Databricks integrations}
\begin{itemize}
    \item Create a service principal in Azure and provide it Azure blob storage role.
    \item Assign the same service principal a role in the Azure Databricks workspace (contributor).
    \item Create tables and schemas that adhere to the data engineering architecture.
    \item Create an API that should be used to push data from the gold layer to an external PostgreSQL.
\end{itemize}

\section{Discussion}

\section{Conclusion}

\section{Appendix}
\subsection{Code}

\end{document}